\section{引言}
\label{sec:intro}

安全团队开始把告警分诊和调查交给 LLM 智能体~\citep{guide,cortex,sentinelrl}。智能体收到一条告警，查询日志，然后报告是否发生了攻击、攻击涉及了什么。这样的智能体能否被信任去关闭告警，很大程度上由基准决定：ExCyTIn-Bench~\citep{excytin}、SIABench~\citep{siabench} 等基准报告智能体答对的比例，而领域把这些数字当作调查能力的度量。调查正是部署智能体的理由。只会复述告警内容的智能体对分析员没有价值，而且恰好会在攻击者刻意制造的情形中失败：恶意活动看起来像日常操作，也没有任何东西点出它的名字。

只有当题目不调查就答不出来时，基准才在衡量调查。以一道 ExCyTIn 题目为例，它由一条三跳路径生成：从一条关于计划任务的告警出发，经过主机，到达一条相关告警，答案是后者记录的 IP 地址。题目问的是"相关的 `Anomaly detected in ASEP registry' 告警中记录的 IP 地址"。按这个告警名做一次 SQL 查询就能拿到答案，生成题目时设想的那次跳转并不需要（\cref{fig:example}）。如果这样的题目很多，智能体不调查也能拿分，而分数对真正重要的情形毫无说明。

基准论文报告模型成绩，却很少报告一个不调查的规则能拿多少分。安全领域有针对学习型检测器做重新评估的传统，问的正是这个问题：简单基线可以与复杂的溯源检测器相当~\citep{bilot2025simpler}，虚假相关会抬高实验室结果~\citep{arp2022dos}。机器学习中同样的问题称为捷径学习~\citep{geirhos2020shortcut,gururangan2018artifacts}。近期有工作为事件响应提出奖励"发现新证据"的新基准~\citep{sirbench}，但已在使用的基准还没有被审计过，在这些基准上评测的智能体是否利用了其中的捷径，也没有被测量过。

\begin{figure}[t]
  \centering
  \includegraphics[width=\columnwidth]{fig_floor.pdf}
  \caption{不含 LLM 的查表在 ExCyTIn-Bench 上超过了报告的 \oModels{} 个智能体中的 \oBelowGraph{} 个。柱：各模型在 589 道测试题上的报告平均奖励~\citep{excytin}。线：在同一批题目上，两种不调查、只查找题目点名告警的规则，按精确匹配计分（阴影：按事件重抽的 95\% 区间）。深色柱高于图查表。}
  \label{fig:floor}
\end{figure}

我们审计了四个面向 LLM 安全运营智能体的公开基准和数据集（ExCyTIn-Bench、SIABench、GUIDE、OTRF Security-Datasets）以及我们自己构建的三个分诊环境。对每一个，我们写一条固定规则：只使用智能体能看到的信息，不含 LLM，并跳过基准设想的调查，再把它与报告的模型成绩比较。我们发现四类捷径：题目点出了答案所在的证据（S1）；标签跟着打标签的组织走（S2）；数据生成过程留下线索，例如攻击网络的地址（S3）；在人工构建的环境中，对证据模型执行固定程序就能解出所有案例（S4）。ExCyTIn 的作者公开了 14 次智能体运行的完整日志，因此我们进一步逐题检验智能体的成功是否来自捷径。每次确认性运行之前，假设都已登记、规则都已冻结。

\paragraph{结果} 每个被审计的基准都存在一条不调查的规则，能达到报告成绩中相当大的一部分。在 ExCyTIn 上，\oNamed\% 的测试题点出了答案所在的告警，查找这条告警可以精确答对 \oGraph\%，高于报告的 \oModels{} 个模型中 \oBelowGraph{} 个的平均奖励（\cref{fig:floor}）。在 SIABench 上，"当且仅当告警涉及 172.16.0.1 时判为真阳性"在全部 35 条告警上都正确。在 GUIDE 上，沿用组织过去对同一检测器的判定，在之后的事件上正确率为 \gHistPct\%。在 OTRF 录制中，仅凭告警文字就能区分 \otN{} 条告警中的 \otRuleCorrect{} 条。然而 ExCyTIn 的智能体并没有使用这条捷径：题目点名答案告警时成功率并不更高（\lgModels{} 次运行中 \lgPos{} 次为正，$p=\lgP$）；只需一次查询时它们也发出同样多的查询；在抗捷径题目上排名几乎不变（Kendall $\tau=\lgTau$）。更深层的问题是，成功率几乎不随题目所需的调查而变化：最强的智能体在 140 道一次查询即可答对的题目中答错 \tcOpusFail{} 道。

\paragraph{贡献}
\begin{itemize}\setlength\itemsep{1pt}
\item 给出调查型基准中"捷径"的定义和四类捷径的分类，每类配一个不调查的基线（\cref{sec:method}）。
\item 审计四个公开基准和三个人工构建的环境，表明每一个都存在能达到或超过报告成绩相当部分的规则（\cref{sec:floors}）。
\item 对 ExCyTIn 公开的 14 次智能体运行做逐题分析，表明智能体没有利用捷径，其成功率也不随所需调查而变化（\cref{sec:scores}）。
\item 给出面向基准构建者的检查项，证明攻击与良性"孪生"配对可以消除单条观测捷径，并公开 ExCyTIn 的逐题捷径标记（\cref{sec:design}）。
\end{itemize}
